\documentclass{leveldoc}
\DocUses{foundations,first_bend}
\DocFingerprint{7e6dc5c9554c57da}

\begin{document}

\section{What the level asks}\label{sec:asks}

The game's text:
\leveldescription

Three shots leave $(\val{xLaunch}, \val{yLaunch})$ with the same particle
($q = \val{q}$, $m = \val{m}$) and the same direction, $+x$. The kinetic energies are
$T = \val{TLow}$, $T = \val{TMid}$ and $T = \val{THigh}$. Each shot has its own
detector, a rectangle of width $\val{xDhi} - \val{xDlo}$ standing just to the left of
$x = \val{xC}$: $\val{yDloLow} \le y \le \val{yDhiLow}$ for the lowest energy,
$\val{yDloMid} \le y \le \val{yDhiMid}$ for the middle one, and
$\val{yDloHigh} \le y \le \val{yDhiHigh}$ for the highest. An entry counts at every
direction and every kinetic energy. The higher energy has the lower box. The player
places at most $\val{maxCharges}$ charges, of either sign, inside the region
$\val{xRegLo} \le x \le \val{xRegHi}$, $\val{yRegLo} \le y \le \val{yRegHi}$, and no
magnets. A flight ends when the particle touches a charge, leaves the world, enters
the detector or runs out of time (\xsecref{foundations}{sec:board});
$t_{\max} = \val{tmax}$.

The reference solution is one attractive charge, $Q = \val{Q}$ at
$(\val{xC}, \val{yC})$, directly below the source. The coupling has magnitude
$\val{kAbs}$. Each energy then meets the right face of its own detector. This
document follows the three orbits in closed form, and the same charge at the two
neighbouring magnitudes $\val{kAbsWeak}$ and $\val{kAbsStrong}$.

\section{Setup}\label{sec:setup}

The level file, field by field (the appendix reproduces it):
\setuptable

In symbols: $c = \val{c}$, the particle has radius $a = \val{aRad}$, and the charge
has radius $R = \val{R}$. The launch radius, from the charge to the source, is
$r_0 = \val{rZero}$. The direction is perpendicular to that radius. For a kinetic
energy $T$,
\[
  \gamma_0 = 1 + \frac{T}{mc^2}, \qquad
  \abs{\vect{p}_0} = \sqrt{\frac{T^2}{c^2} + 2mT}, \qquad
  v_0 = \frac{\abs{\vect{p}_0}}{\gamma_0 m}.
\]
At $T = \val{TMid}$ this is $\gamma_0 = \val{gammaMid}$,
$\abs{\vect{p}_0} = \val{pMid}$ and $v_0 = \val{vMid}$. The world is
$-1 \le x \le \val{xWorld}$, $-1 \le y \le \val{yWorld}$. The reference coupling is
$\kappa = qQ = \val{kappa}$.

\begin{figure}[htbp]
\centering
\begin{tikzpicture}
\begin{axis}[axis equal image, width=0.88\linewidth, xmin=-1.8, xmax=32.4, ymin=-1.8, ymax=21.8,
  xtick={0,5,...,30}, ytick={0,5,...,20}, tick label style={font=\small},
  xlabel={$x$}, ylabel={$y$}, label style={font=\small}]
  \draw[black!45, dashed] (-1,-1) rectangle (31,21);
  \draw[black!50] (0,0) rectangle (30,20);
  \draw[black!40, dashed] (1,3) rectangle (12,16);
  \fill[green!50!black, fill opacity=0.22] (1,4) rectangle (3,5);
  \draw[green!40!black] (1,4) rectangle (3,5);
  \fill[green!50!black, fill opacity=0.22] (1,3) rectangle (3,4);
  \draw[green!40!black] (1,3) rectangle (3,4);
  \fill[green!50!black, fill opacity=0.22] (1,1) rectangle (3,2);
  \draw[green!40!black] (1,1) rectangle (3,2);
  \fill[red!70!black, fill opacity=0.35] (\rawval{xC},\rawval{yC}) circle[radius=0.3];
  \draw[red!70!black] (\rawval{xC},\rawval{yC}) circle[radius=0.3];
  \addplot[very thick, blue!70!black] table {pathlow.dat};
  \addplot[thick, orange!80!black] table {pathmid.dat};
  \addplot[thick, purple!70!black] table {pathhigh.dat};
  \addplot[thin, black!55] table {pathfloor.dat};
  \addplot[thin, red!60!black] table {pathleft.dat};
  \draw[-{Stealth[length=6pt]}, thick] (\rawval{xLaunch},\rawval{yLaunch}) --
    (\rawval{arrowX},\rawval{arrowY});
  \fill (\rawval{xLaunch},\rawval{yLaunch}) circle[radius=0.11];
  \node[font=\scriptsize, red!60!black, anchor=south] at (4,10.45) {$Q$};
  \node[font=\scriptsize, green!30!black, anchor=east] at (0.85,4.5) {low};
  \node[font=\scriptsize, green!30!black, anchor=east] at (0.85,3.5) {mid};
  \node[font=\scriptsize, green!30!black, anchor=east] at (0.85,1.5) {high};
\end{axis}
\end{tikzpicture}
\caption{The reference charge bends three energies into three detectors. Blue is
$T = \val{TLow}$, orange $T = \val{TMid}$, purple $T = \val{THigh}$. Each enters its
box through the right face. Grey is the weaker charge on the highest energy: it meets
the floor. Thin red is the stronger charge on the lowest energy: it meets the left
edge of the world. The disc has radius $R$. The dashed rectangle is the region in
which the player may place charges.}
\label{fig:board}
\end{figure}

\section{Principles and model}\label{sec:model}

\begin{principles}
  \item The field of a charged sphere and the force on a non-polarizable particle,
    \xeqref{foundations}{eq:coulomb} and \xeqref{foundations}{eq:sphereforce}.
  \item The relativistic equation of motion \xeqref{foundations}{eq:eom}, the launch
    \xeqref{foundations}{eq:launch} and the energy \xeqref{foundations}{eq:energy}. The
    motion stays in the plane (\xsecref{foundations}{sec:plane}).
  \item The board's rules (\xsecref{foundations}{sec:board}). The first entry into the
    detector decides. This level adds no direction and no kinetic window.
  \item One fixed centre separates. The orbit is \xeqref{first_bend}{eq:general}, and
    the time of an apsidal launch is \xeqref{first_bend}{eq:time}.
  \item Newton's method \xeqref{foundations}{eq:newton} for the face, and Simpson's
    rule \xeqref{foundations}{eq:simpson} for the radiated power
    \xeqref{foundations}{eq:larmor}, with the acceleration
    \xeqref{first_bend}{eq:acceleration}.
\end{principles}

The level has no magnets, so $\vect{B} = \vect{0}$. Outside the charge the force is
the Coulomb force of $\kappa$ placed at $(\val{xC}, \val{yC})$.

\begin{assumption}{The charge is held in place}\label{as:fixed}
The charge does not recoil. That is a rule of the game
(\xsecref{foundations}{sec:radiation}).
\end{assumption}

\begin{assumption}{No radiation}\label{as:radiation}
The radiated energy is not taken from the particle. Over the nine flights of this
document the largest fraction of the launch energy is $\val{radFracMax}$
(§\ref{sec:hand}). The game makes the same approximation and reports the same integral
(\physref{8}).
\end{assumption}

\begin{assumption}{A point particle}\label{as:point}
$a = 0$, and the particle does not polarize, so \xeqref{foundations}{eq:sphereforce}
is exact. On the reference orbits the closest approach is $\val{rMinLow}$, on the
lowest energy, with clearance $\val{clearLow}$ outside the charge. The highest energy
is closest at the launch, where the clearance is $\val{clearHigh}$. Both are exact
values of the closed orbit (§\ref{sec:orbit}), far above $R$.
\end{assumption}

\begin{assumption}{The face root is exact for this comparison}\label{as:step}
The orbit and the time are closed forms. The face is a root of Newton's method whose
residual is at most $\val{faceResMid}$ (§\ref{sec:hand}), and the time agrees with an
independent quadrature to $\val{tQuadDiff}$. The tolerance in §\ref{sec:comparison} is
$10^{-9}$, ten times the accuracy of a verified flight.
\end{assumption}

\section{Derivation}\label{sec:derivation}

\subsection{The circular condition}\label{sec:circle}

On a circle the force is perpendicular to the velocity, so
\xeqref{first_bend}{eq:acceleration} reduces to $\vect{a} = \vect{F}/(\gamma m)$. The
speed is constant and the path has curvature $1/r$, so $\abs{\vect{a}} = v^2/r$. With
$\abs{\vect{F}} = \abs{\kappa}/r^2$ this is
\begin{equation}\label{eq:circle}
  \abs{\kappa} = \gamma m v^2 r = p\,v\,r,
\end{equation}
since $p = \gamma m v$. The level's condition is \eqref{eq:circle} at the launch
radius $r_0$. It is not the Newtonian $2Tr_0$: that limit replaces $pv$ by $2T$ only
as $v/c \to 0$.

At the three launch energies, $p_0 v_0 r_0$ equals $\gamma_0 m v_0^2 r_0$: the two
writings are $\val{circLow}$ and $\val{circAltLow}$, $\val{circMid}$ and
$\val{circAltMid}$, $\val{circHigh}$ and $\val{circAltHigh}$. The reference magnitude
is $\val{kAbs}$. The differences $p_0 v_0 r_0 - \abs{\kappa}$ are $\val{circGapLow}$,
$\val{circGapMid}$ and $\val{circGapHigh}$. The lowest energy and the middle energy
are more tightly bound than the circle; the highest energy is less tightly bound.
None of the three is the circle. The middle energy is the nearest: its gap is the
small negative $\val{circGapMid}$.

\subsection{What the energy allows}

The potential at the launch is $\kappa/r_0 = \val{VZeroMid}$. Then
$W - mc^2 = T + \kappa/r_0$ is $\val{WbindLow}$, $\val{WbindMid}$ and
$\val{WbindHigh}$. All three are negative. Far from the charge the potential falls to
$0$, which lies above $W - mc^2$, so \xeqref{foundations}{eq:energy} does not allow
the particle to reach infinity. The orbit stays at a finite radius, and the
integration has to decide which boundary it meets.

\subsection{The orbit}\label{sec:orbit}

The launch radius is vertical and the momentum is horizontal, so
$p_{r,0} = \val{prZero}$ and $u'_0 = \val{duZero}$. The launch is an apsis. The
angular momentum is negative, $L_z = -r_0\abs{\vect{p}_0}$, so $\sigma = \val{sigma}$:
the particle turns clockwise from the top, and $\theta$ is the angle swept from the
launch. \xeqref{first_bend}{eq:general} at $u'_0 = 0$ is
\[
  u(\theta) = B + A\cos\Gamma\theta, \qquad A = u_0 - B,
\]
with $B$ and $\Gamma^2$ as in \xeqref{first_bend}{eq:orbit} and
\xeqref{first_bend}{eq:general}. Here $\Gamma^2 > 0$. The positive root of
\xeqref{first_bend}{eq:amplitude} equals $\abs{A}$; the sign of $A$ is the sign of
$u_0 - B$.

$A < 0$ means the launch is the outer apsis: $u$ increases and the orbit falls
inward. That is the lowest energy, $A = \val{ALow}$, and the middle energy,
$A = \val{AMid}$. $A > 0$ means the launch is the inner apsis and the orbit swings
outward. That is the highest energy, $A = \val{AHigh}$. On every flight of this
document $B > \abs{A}$, so $u$ stays positive and the particle never reaches
infinity, in agreement with the sign of $W - mc^2$.

$\Gamma$ is less than $1$ by $\val{oneGLow}$, $\val{oneGMid}$ and $\val{oneGHigh}$.
The relativistic orbit is this apsidal form, not a fixed ellipse: the radial period
and a full turn differ by that fraction. Over the arc to the detector the whole
effect is of that size.

The position is the polar angle $\varphi_0 + \sigma\theta$ about
$(\val{xC}, \val{yC})$, with $\varphi_0 = \val{phiZero}$. Because the charge is
directly below the source, $\varphi_0$ is a right angle, and
\[
  x(\theta) = \val{xC} + \frac{\sin\theta}{u(\theta)}, \qquad
  y(\theta) = \val{yC} + \frac{\cos\theta}{u(\theta)}.
\]
The bottom of the orbit is $\theta = \pi$. The detectors sit at $x = \val{xDhi}$, one
cell to the left of the charge, so a clockwise particle meets that line only after
passing the bottom.

\subsection{The time}\label{sec:time}

\xeqref{first_bend}{eq:dtdth} and the change of variable $s = \Gamma\theta$ give
\xeqref{first_bend}{eq:time},
\[
  t(\theta) = \frac{W\,I_2(\Gamma\theta) - \kappa\,I_1(\Gamma\theta)}{c^2 L\,\Gamma},
  \qquad I_n(s) = \int_0^s \frac{\dd s'}{(A\cos s' + B)^n}.
\]
\xeqref{first_bend}{eq:i2} still gives $I_2$ from $I_1$. Its derivation is a
differentiation and uses $A \ne 0$, which holds here: the smallest $\abs{A}$ is
$\val{AMid}$, on the energy nearest the circle. \xeqref{first_bend}{eq:i1} does not
carry over. That antiderivative is an $\operatorname{artanh}$ because $A > \abs{B}$ on
the repulsive orbit of Level~1. Here $B > \abs{A}$, so the same half-angle
substitution produces an arctangent. With $\Delta = \sqrt{B^2 - A^2}$ and
$0 \le s < \pi$,
\begin{equation}\label{eq:i1attr}
  I_1(s) = \frac{2}{\Delta}
  \arctan\!\left(\tan\frac{s}{2}\sqrt{\frac{B - A}{B + A}}\right).
\end{equation}
As $s \to \pi$ this tends to $\pi/\Delta$. For $\pi \le s < 2\pi$ the angle
$s - \pi$ runs through a cosine of the opposite sign, $\cos(s'+\pi) = -\cos s'$, and
\[
  I_1(s) = \frac{\pi}{\Delta} + \frac{2}{\Delta}
  \arctan\!\left(\tan\frac{s - \pi}{2}\sqrt{\frac{B + A}{B - A}}\right).
\]
Every flight of this document has $s < 2\pi$, which is checked. The three reference
values $\val{sLow}$, $\val{sMid}$ and $\val{sHigh}$ all exceed $\pi$, so those times
use the split.

\subsection{Which face, and where}\label{sec:event}

The first event is read from chords of angle $\val{dth}$. On these nine flights the
greatest distance from a chord to the orbit at the chord's midpoint is
$\val{sagMax}$, far below the one-cell height of a detector, so a chord cannot skip a
face or reorder two events. A tie would keep the board's order
(\xsecref{foundations}{sec:board}): the charge, then the world, then the detector.
The winning face is then solved by Newton's method \xeqref{foundations}{eq:newton} on
$f(\theta) = x(\theta) - x_\mathrm{f}$ or $f(\theta) = y(\theta) - y_\mathrm{f}$, with
the derivative of §\ref{sec:orbit}. The reported face coordinate is the face value.
The residual of $f$, before that assignment, is the check that the root lies on the
face. The momentum is \xeqref{first_bend}{eq:momentum} at the same $\theta$, and the
time is §\ref{sec:time}.

The radiated energy is $\int P\,(\dd t/\dd\theta)\,\dd\theta$ from the launch to that
root. The integrand is evaluated in binary64 and summed by Simpson's rule with
$\val{nSimp}$ panels. Halving the panel count estimates the relative error as
$(S_n - S_{n/2})/15$.

\section{The reference solution by hand}\label{sec:hand}

\begin{procedure}
  \item\label{st:launch} \emph{The launches} (\xeqref{foundations}{eq:launch}).
    The lowest energy has $\gamma_0 = \val{gammaLow}$ and
    $\abs{\vect{p}_0} = \val{pLow}$. The middle energy has
    $\gamma_0 = \val{gammaMid}$ and $\abs{\vect{p}_0} = \val{pMid}$. The highest has
    $\gamma_0 = \val{gammaHigh}$ and $\abs{\vect{p}_0} = \val{pHigh}$. All three are
    aimed along $+x$, and each is an apsis.
  \item\label{st:circle} \emph{The circle} \eqref{eq:circle}. Against
    $\abs{\kappa} = \val{kAbs}$ the three values of $p_0 v_0 r_0$ leave the gaps
    $\val{circGapLow}$, $\val{circGapMid}$ and $\val{circGapHigh}$.
  \item\label{st:orbit} \emph{The constants.} $W - mc^2$ is $\val{WbindLow}$,
    $\val{WbindMid}$ and $\val{WbindHigh}$, all negative. The orbit constants are
    $L = \val{LLow}$, $\Gamma = \val{GLow}$, $A = \val{ALow}$, $B = \val{BLow}$ for
    the lowest energy; $L = \val{LMid}$, $\Gamma = \val{GMid}$, $A = \val{AMid}$,
    $B = \val{BMid}$ for the middle; and $L = \val{LHigh}$, $\Gamma = \val{GHigh}$,
    $A = \val{AHigh}$, $B = \val{BHigh}$ for the highest. The amplitude check differs
    from $\abs{A}$ by $\val{AdiffLow}$, $\val{AdiffMid}$ and $\val{AdiffHigh}$.
  \item\label{st:newton} \emph{The middle energy, on $x = \val{xDhi}$.} Newton's
    method from the chord that first meets the face:
    \begin{center}\small\val{newton}\end{center}
    Four iterations leave the residual $\val{faceResMid}$. The root is
    $\theta_\mathrm{e} = \val{thMid}$, which is $\val{pastMid}^\circ$ past the bottom,
    and $s_\mathrm{e} = \val{sMid}$. Then $I_1 = \val{IOne}$, $I_2 = \val{ITwo}$ and
    $t_\mathrm{e} = \val{RefMidT}$. The same integral by quadrature differs by
    $\val{tQuadDiff}$. The entry is $(\val{RefMidX}, \val{RefMidY})$, so
    $y_\mathrm{e}$ lies between $\val{yDloMid}$ and $\val{yDhiMid}$. The momentum is
    $(\val{RefMidPx}, \val{RefMidPy})$, at $\val{angMid}^\circ$ to $+x$. Its radial
    and angular parts are $\val{prE}$ and $\val{pphiE}$.
  \item\label{st:entries} \emph{The other two energies.} The lowest meets the same
    face at $\theta = \val{thLow}$, $\val{pastLow}^\circ$ past the bottom, at
    $(\val{RefLowX}, \val{RefLowY})$ and $t = \val{RefLowT}$, with momentum
    $(\val{RefLowPx}, \val{RefLowPy})$, at $\val{angLow}^\circ$. The ordinate lies
    between $\val{yDloLow}$ and $\val{yDhiLow}$.

    The highest meets it at $\theta = \val{thHigh}$, $\val{pastHigh}^\circ$ past the
    bottom, at $(\val{RefHighX}, \val{RefHighY})$ and $t = \val{RefHighT}$, with
    momentum $(\val{RefHighPx}, \val{RefHighPy})$, at $\val{angHigh}^\circ$. The
    ordinate lies between $\val{yDloHigh}$ and $\val{yDhiHigh}$.

    The three residuals are $\val{faceResLow}$, $\val{faceResMid}$ and
    $\val{faceResHigh}$. The low-energy residual is zero.
  \item\label{st:check} \emph{The checks.} The arrival kinetic energies are
    $\val{TarrLow}$, $\val{TarrMid}$ and $\val{TarrHigh}$. The two that fall inward
    arrive above their launch energy; the one that swings outward arrives below it.
    The closest approaches are $\val{rMinLow}$, $\val{rMinMid}$ and $\val{rMinHigh}$,
    each well above $R$. Every reference time matches its quadrature to the tolerance
    of Step~\ref{st:newton}.
  \item\label{st:rad} \emph{The radiation.} Simpson's rule as in §\ref{sec:event}
    gives $\val{RefLowRad}$, $\val{RefMidRad}$ and $\val{RefHighRad}$. The relative
    step-halving uncertainties are $\val{RefLowRadErr}$, $\val{RefMidRadErr}$ and
    $\val{RefHighRadErr}$.
\end{procedure}

\begin{result}[Three energies, three boxes]
The lowest energy enters its detector at $(\val{RefLowX}, \val{RefLowY})$ at
$t = \val{RefLowT}$. The middle energy enters at $(\val{RefMidX}, \val{RefMidY})$ at
$t = \val{RefMidT}$. The highest energy enters at $(\val{RefHighX}, \val{RefHighY})$
at $t = \val{RefHighT}$. Each ordinate lies in its own box, and each arrival is past
the bottom of the orbit. The level is solved.
\end{result}

The arithmetic is the closed form of §\ref{sec:orbit} and §\ref{sec:time}. The only
iteration is Step~\ref{st:newton}, and the same root on the other two energies. The
radiated energy is the one numerical integral, used for \ref{as:radiation}.

\section{Comparison with the game}\label{sec:comparison}

\comparisontable

The nine outcomes agree, and so do the nine end positions and times. The reference
momenta are compared as well. The tolerance on a state is $\val{tolFlight}$. The
three reference radiated energies are compared relatively, with tolerance
$\val{tolRad}$. Ten times the largest Simpson uncertainty $\val{radErrMax}$ lies
below $10^{-4}$, so the radiation tolerance is that floor. The state tolerance and
the radiation tolerance were fixed from these estimates before the comparison.

The largest absolute difference is the time of the weaker charge's highest energy,
$\val{cmpWeakHighTDiff}$. Its free coordinate, the abscissa at the floor, differs by
$\val{cmpWeakHighXDiff}$. The document's root and time are closed forms: the time
matches a quadrature to $\val{tQuadDiff}$, and the face residual is at most
$\val{faceResMid}$. The difference is the size of the game's verified accuracy, about
$10^{-10}$ (\xsecref{foundations}{sec:verification}).

The quadrature uncertainties of the three reference flights are
$\val{RefLowRadErr}$, $\val{RefMidRadErr}$ and $\val{RefHighRadErr}$, inside the
relative tolerance. The largest relative difference is the highest energy,
$\val{cmpRefHighRadDiff}$.

\section{Other solutions}\label{sec:others}

The same point with either neighbouring magnitude fails every shot. Both flights are
the orbit of §\ref{sec:orbit} with the face of §\ref{sec:event}.
\begin{enumerate}[label=\arabic*.]
  \item $Q = \val{Qweak}$, so $\abs{\kappa} = \val{kAbsWeak}$, which is below every
    $p_0 v_0 r_0$. All three orbits swing outward. The lowest and the middle energy
    pass the line $x = \val{xDhi}$ underneath their detectors and leave through the
    left edge of the world. The highest energy meets the floor first.

    The lowest energy crosses $x = \val{xDhi}$ at $y = \val{WeakLowCrossY}$, below
    $\val{yDloLow}$, and leaves at $(\val{WeakLowX}, \val{WeakLowY})$, at
    $t = \val{WeakLowT}$.

    The momentum there is $(\val{WeakLowPx}, \val{WeakLowPy})$, at
    $\val{WeakLowAng}^\circ$. The radiated energy is $\val{WeakLowRad}$.

    The middle energy crosses $x = \val{xDhi}$ at $y = \val{WeakMidCrossY}$, below
    $\val{yDloMid}$, and leaves at $(\val{WeakMidX}, \val{WeakMidY})$, at
    $t = \val{WeakMidT}$.

    The momentum is $(\val{WeakMidPx}, \val{WeakMidPy})$, at $\val{WeakMidAng}^\circ$.
    The radiated energy is $\val{WeakMidRad}$.

    The highest energy meets the bottom face of the world at
    $(\val{WeakHighX}, \val{WeakHighY})$, at $t = \val{WeakHighT}$.

    The momentum is $(\val{WeakHighPx}, \val{WeakHighPy})$, at
    $\val{WeakHighAng}^\circ$, still outward of the charge. The radiated energy is
    $\val{WeakHighRad}$.
  \item $Q = \val{Qstrong}$, so $\abs{\kappa} = \val{kAbsStrong}$, which is above
    every $p_0 v_0 r_0$. All three orbits fall inward, cross $x = \val{xDhi}$ above
    their detectors, and leave through the left edge of the world.

    The lowest energy crosses at $y = \val{StrongLowCrossY}$, above $\val{yDhiLow}$,
    and leaves at $(\val{StrongLowX}, \val{StrongLowY})$, at $t = \val{StrongLowT}$.

    The momentum is $(\val{StrongLowPx}, \val{StrongLowPy})$, at
    $\val{StrongLowAng}^\circ$. The radiated energy is $\val{StrongLowRad}$.

    The middle energy crosses at $y = \val{StrongMidCrossY}$, above $\val{yDhiMid}$,
    and leaves at $(\val{StrongMidX}, \val{StrongMidY})$, at $t = \val{StrongMidT}$.

    The momentum is $(\val{StrongMidPx}, \val{StrongMidPy})$, at
    $\val{StrongMidAng}^\circ$. The radiated energy is $\val{StrongMidRad}$.

    The highest energy crosses at $y = \val{StrongHighCrossY}$, above
    $\val{yDhiHigh}$, and leaves at $(\val{StrongHighX}, \val{StrongHighY})$, at
    $t = \val{StrongHighT}$.

    The momentum is $(\val{StrongHighPx}, \val{StrongHighPy})$, at
    $\val{StrongHighAng}^\circ$. The radiated energy is $\val{StrongHighRad}$.
\end{enumerate}

The reference magnitude is the one that places each crossing of $x = \val{xDhi}$
inside its own box. One step weaker, the crossings fall below the boxes, and the
highest energy meets the floor before that line. One step stronger, the crossings
rise above the boxes. Section~\ref{sec:comparison} records the outcomes, and
Figure~\ref{fig:board} draws the two misses that leave through the floor and through
the left edge.

\section{What the game computes}\label{sec:engine}

The game integrates \xeqref{foundations}{eq:eom} with its eighth-order Runge--Kutta
method at the tolerances $10^{-10}$ and $10^{-12}$ (\physref{5}). It finds the crossing
on its interpolant (\physref{6}) and takes the first entry into the box as the outcome
(\xsecref{foundations}{sec:board}). It verifies that outcome by comparing its two runs
(\physref{7}). The model is this document's, including \ref{as:radiation} (\physref{8}).
The agreement in §\ref{sec:comparison} is between the closed orbit of §\ref{sec:orbit}
and that integration.

\section{Further reading}\label{sec:reading}

\begin{itemize}
  \item \levelref{first_bend}: the orbit \xeqref{first_bend}{eq:general}, the apsidal
    time \xeqref{first_bend}{eq:time}, and the momentum at a point of the orbit. The
    integral $I_1$ there is an $\operatorname{artanh}$; \eqref{eq:i1attr} is the same
    integral for $B > \abs{A}$.
  \item \xsecref{foundations}{sec:numerics}: Newton's method and Simpson's rule.
  \item \xsecref{foundations}{sec:board}: the order of events.
\end{itemize}

\levelfileappendix

\end{document}
